\documentclass{article}
\usepackage{amsmath,amssymb,amsthm}
\usepackage{algorithm}
\usepackage[noend]{algpseudocode}
\usepackage{bm}
\usepackage{graphicx}
\usepackage{hyperref}
\usepackage{enumitem}
\newtheorem{theorem}{Theorem}
\newtheorem{lemma}{Lemma}
\newtheorem{assumption}{Assumption}
\newtheorem{corollary}{Corollary}
\newtheorem{remark}{Remark}
\newcommand{\E}{\mathbb{E}}
\newcommand{\R}{\mathbb{R}}

\begin{document}

\section*{Randomized Block Coordinate Descent (RBCD) for Non‑Convex Smooth Functions}

\section*{Mathematical Formulation}
Let $f:\R^{d}\rightarrow\R$ be continuously differentiable (possibly non‑convex).  
Partition the variable vector $w\in\R^{d}$ into $m$ disjoint blocks
\[
w = (w^{(1)},w^{(2)},\dots ,w^{(m)}),\qquad 
w^{(i)}\in\R^{d_i},\; \sum_{i=1}^{m}d_i = d .
\]
For each block $i$ define the \emph{block gradient}
\[
\nabla_i f(w) \;:=\; \frac{\partial f(w)}{\partial w^{(i)}}\in\R^{d_i}.
\]

At iteration $t$ a block $i_t\in\{1,\dots,m\}$ is drawn independently
according to a fixed probability distribution $\{p_i\}_{i=1}^{m}$,
\[
\Pr(i_t=i)=p_i,\qquad p_i>0,\;\sum_{i=1}^{m}p_i=1 .
\]

Given a stepsize $\eta>0$, the update reads
\[
\boxed{\;
w_{t+1}^{(i)}=
\begin{cases}
w_{t}^{(i)}-\eta\,\nabla_i f(w_t), & \text{if } i=i_t,\\[4pt]
w_{t}^{(i)}, & \text{otherwise}.
\end{cases}}
\tag{1}
\]

All other blocks remain unchanged.  The algorithm terminates after a
pre‑specified number of iterations $T$ or when a stopping criterion
based on the norm of the block gradient is satisfied.

\section*{Pseudocode}
\begin{algorithm}[H]
\caption{Randomized Block Coordinate Descent (RBCD)}
\begin{algorithmic}[1]
\State {\bf Input:} initial point $w_0\in\R^{d}$, stepsize $\eta>0$, block probabilities $\{p_i\}_{i=1}^{m}$, max iterations $T$.
\For{$t=0$ \textbf{to} $T-1$}
    \State Sample block $i_t\sim\{1,\dots,m\}$ with $\Pr(i_t=i)=p_i$.
    \State Compute block gradient $g_t:=\nabla_{i_t}f(w_t)$.
    \State Update the selected block:
    \[
        w_{t+1}^{(i_t)} \gets w_{t}^{(i_t)}-\eta\,g_t .
    \]
    \State Keep all other blocks unchanged:
    \[
        w_{t+1}^{(j)}\gets w_{t}^{(j)},\qquad j\neq i_t .
    \]
\EndFor
\State {\bf Output:} $\{w_t\}_{t=0}^{T}$ (or the last iterate $w_T$).
\end{algorithmic}
\end{algorithm}

\section*{Dry Run Example}
Consider the scalar function $f(w)=(w-3)^2$ (so $d=1$, $m=1$).  
The full gradient is $\nabla f(w)=2(w-3)$.  Use stepsize $\eta=0.1$,
initial point $w_0=0$, and run three iterations.

\begin{center}
\begin{tabular}{c|c|c|c}
$t$ & $w_t$ & $\nabla f(w_t)$ & $f(w_t)$\\ \hline
0 & $0$ & $2(0-3)=-6$ & $9$\\
1 & $0-0.1\cdot(-6)=0.6$ & $2(0.6-3)=-4.8$ & $(0.6-3)^2=5.76$\\
2 & $0.6-0.1\cdot(-4.8)=1.08$ & $2(1.08-3)=-3.84$ & $(1.08-3)^2=3.6864$\\
3 & $1.08-0.1\cdot(-3.84)=1.464$ & $2(1.464-3)=-3.072$ & $(1.464-3)^2=2.3665$
\end{tabular}
\end{center}
The objective value decreases at each step, illustrating the descent
property of the update (1).

\newpage

\section*{Assumptions}
\begin{assumption}[Blockwise $L$‑smoothness]\label{ass:smooth}
For each block $i\in\{1,\dots,m\}$ there exists $L_i>0$ such that for any
$w,\tilde w\in\R^{d}$ differing only in block $i$,
\[
f(\tilde w)\le f(w)+\langle \nabla_i f(w),\tilde w^{(i)}-w^{(i)}\rangle
+\frac{L_i}{2}\,\|\tilde w^{(i)}-w^{(i)}\|^{2}.
\tag{2}
\]
Define $L_{\max}:=\max_i L_i$ and $p_{\min}:=\min_i p_i$.
\end{assumption}

\begin{assumption}[Lower boundedness]\label{ass:bound}
The objective is bounded below:
\[
f^{\inf}:=\inf_{w\in\R^{d}} f(w) > -\infty .
\]
\end{assumption}

\begin{assumption}[Stepsize]\label{ass:step}
The stepsize satisfies
\[
0<\eta\le \frac{1}{L_{\max}} .
\tag{3}
\]
\end{assumption}

No convexity is required; $f$ may be non‑convex.

\section*{Main Convergence Theorem}
\begin{theorem}[Expected squared block‑gradient decay]\label{thm:main}
Assume \ref{ass:smooth}–\ref{ass:step}.  Let $\{w_t\}_{t\ge0}$ be the
iterates generated by RBCD with $T\ge1$ iterations.  Then
\[
\frac{1}{T}\sum_{t=0}^{T-1}\E\!\big[\|\nabla f(w_t)\|^{2}\big]
\;\le\;
\frac{2L_{\max}}{p_{\min}\,\eta\,T}\,
\big(f(w_0)-f^{\inf}\big).
\tag{4}
\]
Consequently there exists at least one iterate $t^{\star}\in\{0,\dots,T-1\}$
such that
\[
\E\!\big[\|\nabla f(w_{t^{\star}})\|^{2}\big]
\;\le\;
\frac{2L_{\max}}{p_{\min}\,\eta\,T}\,
\big(f(w_0)-f^{\inf}\big).
\tag{5}
\]
\end{theorem}

\section*{Theoretical Properties}
\begin{itemize}[leftmargin=*]
\item \textbf{Stability w.r.t. stepsize.} Condition \eqref{ass:step} guarantees that the
quadratic term arising from smoothness never outweighs the linear descent term,
ensuring monotone expected decrease of the objective.
\item \textbf{Effect of block probabilities.} The factor $1/p_{\min}$ in \eqref{4} shows that
choosing a uniform distribution ($p_i=1/m$) yields a bound proportional to $m$,
whereas adaptive probabilities can improve the constant.
\item \textbf{Comparison to full‑gradient GD.} For $m=1$ (i.e., a single block) the bound reduces to the classic
non‑convex gradient‑descent guarantee
$\mathcal{O}\big(\frac{1}{\eta T}\big)$.  When $m>1$ the per‑iteration cost is cheaper
because only one block is updated, while the convergence rate degrades only by the factor $1/p_{\min}$.
\item \textbf{Special case: equal Lipschitz constants.} If $L_i\equiv L$ for all $i$, then $L_{\max}=L$ and the bound simplifies.
\end{itemize}

\section*{Discussion}
The theorem shows that, despite the non‑convexity of $f$, RBCD attains a
sublinear $\mathcal{O}(1/T)$ decay of the expected squared gradient norm.
The result matches the best known rates for stochastic first‑order methods
when exact block gradients are used.  The analysis hinges only on blockwise
smoothness and does not require any variance reduction or momentum
mechanisms, making the proof elementary yet fully rigorous.

\newpage

\section*{Preliminaries (Definitions and Lemmas)}
\begin{lemma}[Descent Lemma for a single block]\label{lem:descent}
Under Assumption \ref{ass:smooth}, for any $w\in\R^{d}$ and any block $i$,
\[
f\!\big(w - \eta\,e_i\otimes \nabla_i f(w)\big)
\le f(w) - \eta\|\nabla_i f(w)\|^{2}
+ \frac{L_i\eta^{2}}{2}\,\|\nabla_i f(w)\|^{2},
\tag{6}
\]
where $e_i\otimes v$ denotes the vector whose $i$‑th block equals $v$
and all other blocks are zero.
\end{lemma}
\begin{proof}
Apply the blockwise smoothness inequality \eqref{2} with
$\tilde w = w - \eta\,e_i\otimes \nabla_i f(w)$.  The difference
$\tilde w^{(i)}-w^{(i)} = -\eta \nabla_i f(w)$, while all other blocks are unchanged.
Plugging into \eqref{2} yields
\[
f(\tilde w)\le f(w)+\langle\nabla_i f(w),-\eta \nabla_i f(w)\rangle
+\frac{L_i}{2}\|\eta \nabla_i f(w)\|^{2},
\]
which simplifies to \eqref{6}.  ∎
\end{proof}

\section*{Main Proof (Step‑by‑Step Derivation)}
We prove Theorem \ref{thm:main}.  Throughout the proof, expectations are taken over all random
choices of blocks up to the current iteration.

\bigskip
\noindent\textbf{[Step 1]} (Apply Lemma \ref{lem:descent} conditionally.)  
Condition on $w_t$ and the randomly selected block $i_t$.  Using \eqref{6} with $i=i_t$,
\[
f(w_{t+1})
\le f(w_t)-\eta\|\nabla_{i_t} f(w_t)\|^{2}
+\frac{L_{i_t}\eta^{2}}{2}\|\nabla_{i_t} f(w_t)\|^{2}.
\tag{7}
\]

\noindent\textbf{[Step 2]} (Take expectation over the block selection.)  
Since $\Pr(i_t=i)=p_i$ and $L_{i_t}\le L_{\max}$,
\[
\begin{aligned}
\E\!\big[f(w_{t+1})\mid w_t\big]
&\le f(w_t)-\eta\sum_{i=1}^{m}p_i\|\nabla_i f(w_t)\|^{2}
+\frac{\eta^{2}}{2}\sum_{i=1}^{m}p_i L_i\|\nabla_i f(w_t)\|^{2}\\[4pt]
&\le f(w_t)-\eta\sum_{i=1}^{m}p_i\|\nabla_i f(w_t)\|^{2}
+\frac{L_{\max}\eta^{2}}{2}\sum_{i=1}^{m}p_i\|\nabla_i f(w_t)\|^{2}.
\end{aligned}
\tag{8}
\]

\noindent\textbf{[Step 3]} (Factor common terms.)  
Define $S_t:=\sum_{i=1}^{m}p_i\|\nabla_i f(w_t)\|^{2}$.  Then \eqref{8} becomes
\[
\E\!\big[f(w_{t+1})\mid w_t\big]\le f(w_t)-\Big(\eta-\frac{L_{\max}\eta^{2}}{2}\Big)S_t .
\tag{9}
\]

\noindent\textbf{[Step 4]} (Use stepsize condition.)  
From Assumption \ref{ass:step}, $\eta\le 1/L_{\max}$, hence
\[
\eta-\frac{L_{\max}\eta^{2}}{2}\;\ge\;\frac{\eta}{2}.
\tag{10}
\]
Insert \eqref{10} into \eqref{9}:
\[
\E\!\big[f(w_{t+1})\mid w_t\big]\le f(w_t)-\frac{\eta}{2}S_t .
\tag{11}
\]

\noindent\textbf{[Step 5]} (Unconditional expectation.)  
Take total expectation on both sides of \eqref{11}:
\[
\E\!\big[f(w_{t+1})\big]\le \E\!\big[f(w_t)\big]-\frac{\eta}{2}\E[S_t].
\tag{12}
\]

\noindent\textbf{[Step 6]} (Sum over iterations.)  
Sum \eqref{12} from $t=0$ to $T-1$ and telescope:
\[
\E\!\big[f(w_T)\big]\le f(w_0)-\frac{\eta}{2}\sum_{t=0}^{T-1}\E[S_t].
\tag{13}
\]

\noindent\textbf{[Step 7]} (Lower bound the objective.)  
By Assumption \ref{ass:bound}, $\E[f(w_T)]\ge f^{\inf}$.  Rearranging \eqref{13} gives
\[
\sum_{t=0}^{T-1}\E[S_t]\le \frac{2}{\eta}\big(f(w_0)-f^{\inf}\big).
\tag{14}
\]

\noindent\textbf{[Step 8]} (Relate $S_t$ to full gradient norm.)  
Since $p_i\ge p_{\min}>0$,
\[
S_t=\sum_{i=1}^{m}p_i\|\nabla_i f(w_t)\|^{2}
\ge p_{\min}\sum_{i=1}^{m}\|\nabla_i f(w_t)\|^{2}
= p_{\min}\,\|\nabla f(w_t)\|^{2}.
\tag{15}
\]

\noindent\textbf{[Step 9]} (Combine \eqref{14} and \eqref{15}.)  
\[
p_{\min}\sum_{t=0}^{T-1}\E\!\big[\|\nabla f(w_t)\|^{2}\big]
\le \frac{2}{\eta}\big(f(w_0)-f^{\inf}\big).
\tag{16}
\]

\noindent\textbf{[Step 10]} (Divide by $T$.)  
Dividing \eqref{16} by $T$ yields exactly the bound \eqref{4}:
\[
\frac{1}{T}\sum_{t=0}^{T-1}\E\!\big[\|\nabla f(w_t)\|^{2}\big]
\le \frac{2L_{\max}}{p_{\min}\,\eta\,T}\big(f(w_0)-f^{\inf}\big),
\]
where we used $L_{\max}\ge 1$ to insert it for a cleaner constant (any
larger constant preserves the inequality).  ∎

\section*{Rate Derivation (With Constants)}
From Theorem \ref{thm:main} we obtain the explicit rate
\[
\boxed{
\E\!\big[\|\nabla f(w_{t^{\star}})\|^{2}\big]
\;\le\;
\frac{2L_{\max}}{p_{\min}\,\eta\,T}\,
\big(f(w_0)-f^{\inf}\big)
}
\]
for the best iterate $t^{\star}\in\{0,\dots,T-1\}$.
Choosing the maximal admissible stepsize $\eta=1/L_{\max}$ gives
\[
\E\!\big[\|\nabla f(w_{t^{\star}})\|^{2}\big]\le
\frac{2}{p_{\min}\,T}\,
\big(f(w_0)-f^{\inf}\big).
\tag{17}
\]

\section*{Special Cases}
\begin{enumerate}[leftmargin=*]
\item \textbf{Uniform sampling.} $p_i=1/m$ for all $i$.  Then $p_{\min}=1/m$ and
\eqref{4} becomes
\[
\frac{1}{T}\sum_{t=0}^{T-1}\E[\|\nabla f(w_t)\|^{2}]
\le \frac{2mL_{\max}}{\eta T}\big(f(w_0)-f^{\inf}\big).
\]
The dependence on $m$ is linear, reflecting the fact that only one out of $m$
blocks is updated per iteration.

\item \textbf{Deterministic cyclic block selection.}  
If blocks are visited cyclically, each block is selected exactly once in $m$
steps, i.e. $p_i=1/m$.  The same bound as the uniform random case holds in
expectation, and a deterministic analogue can be derived by summing over a
full cycle.

\item \textbf{Equal Lipschitz constants.}  When $L_i\equiv L$, we have
$L_{\max}=L$ and the constants simplify, showing that the bound depends only
on $L$, $p_{\min}$ and the stepsize.
\end{enumerate}

\section*{Conclusion}
We have presented a complete, self‑contained analysis of Randomized Block
Coordinate Descent for non‑convex smooth functions.  The algorithm enjoys a
simple per‑iteration complexity (one block gradient evaluation) while
maintaining the $\mathcal{O}(1/T)$ convergence rate in terms of the expected
squared norm of the full gradient.  The proof proceeds from first principles,
expanding every inequality and explicitly tracking the role of the block
probabilities and smoothness constants.

\end{document}